% This file is the inbox for the lecture in progress. Type today's raw notes straight into
% it: it is \input from the chapter file, so the Overleaf preview builds them live. Once the
% material has been placed into the sections, /integrate empties the file again and leaves
% this header behind. Do not delete the file or its \input line.

\section{Relativisation} % [CLAUDE] as usual, figure out where to put this.

Relativisation is essentially like satisfaction, just for classes. Our ultimate goal is the following: given $(A, R)$ a class structure and a formula $\psi$, we'd want to define $\psi^{(A, R)}$ to be the relativisation (think ``interpretation'') of $\psi$ to $(A, R)$.

Here are some complications with this technique:
\begin{itemize}
    \item $R$ might not be $\in$
    \item There might be parameters involved in the definitions of $A$ and $R$. For instance, if $A$ is a set, then ``$A = \setst{x}{x \in A}$'' takes in a set $A$ as a parameter.
\end{itemize}

Look at $A = \setst{x}{\varphi(x)}$. As a scheme, define $\psi^A\of{\bar{v}}$ (with $\bar{v}$ denoting free variables) by recursion on the complexity of $\psi\of{\bar{v}}$ as follows. We will have atomic, propositional and quantifier steps. Before each step, if $u$ is any variable that occurs in $\psi$, use some other $u$ instead: ie, $A = \setst{x}{\varphi'(x)}$ where $\varphi'$ is $\varphi\of{u / u'}$. % [CLAUDE] make this more precize and add an explanation.
\begin{itemize}
    \item \underline{Atomic:} We interpret equality and membership symbols as equality and membership. Ie,
    \begin{align*}
        \parenth{v_0 = v_1}^A &\text{ is } v_0 = v_1 \\
        \parenth{v_0 \in v_1}^A &\text{ is } v_0 \in v_1 
    \end{align*}

    \item \underline{Propositional:} If $\psi$ is a propositional combination of $\psi_0$ and $\psi_1$, then $\psi^A$ is the same propositional combination of $\psi_0^A$ and $\psi_1^A$.

    \item \underline{Quantifiers:} \sorry % [CLAUDE] just write it out, briefly, in the obvious way
\end{itemize}

So for instance, if $R = \setst{(x, y)}{\varphi_R\of{x, y, q}}$, we'd recursively replace occurrences of $v_0 \in v_1$ in $\psi$ with $\varphi_R\of{v_0, v_1, q}$. % [CLAUDE] wut?

Why are we doing all of this? Here's one reason.

\begin{boxtheorem}[working in $\ZFwoF$]
    For each axiom $\sigma \in \ZF$ (\textbf{including} foundation), $\sigma^{\WF}$.
\end{boxtheorem}
The above reads, ``The relativisation of each $\ZF$ axiom to $\WF$ is true, working in $\ZFwoF$.'' Another equivalent formulation is that for all axioms $\sigma$ of $\ZF$,
\begin{align}
    \ZFwoF \vdash \sigma^{\WF}
\end{align}

Theorem \sorry % [CLAUDE] add cross-ref
yields the following corollary.

\begin{boxcorollary}
    If $\ZFwoF$ is consistent, then so is $\ZF$.
\end{boxcorollary}
\begin{proof}
    We argue by contraposition. Suppose $\ZF$ is inconsistent. Say $\sigma_1, \ldots, \sigma_n$ are axioms of $\ZFwoF$ and
    \begin{align*}
        \set{\sigma_1, \ldots, \sigma_n} \cup \set{F} \vdash \varphi \land \neg\varphi
    \end{align*}
    As a scheme (in the meta-theory), we can show that if $A$ is a non-empty class (say, with no parameters), then whenever $\sigma \vdash \tau$, $\sigma^A \vdash \tau^A$. % [CLAUDE] factor this out as a separate lemma and add 2-3 short sentences of intuition preceding it that tie it to the stuff we did earlier on in the lecture. Replace "Apply this" below with a cross-reference after you've factored this out.

    Apply this to what we have to get that
    \begin{align*}
        \set{\sigma_1^{\WF}, \ldots, \sigma_n^{\WF}} \cup \set{F^{\WF}} \vdash \varphi^{\WF} \land \neg \varphi^{\WF}
    \end{align*}
\end{proof}

Together with the above theorem scheme, we get that $\ZFwoF$ is inconsistent. % [CLAUDE] justify

\subsection{Complexity of Formulae}

% [TODO] 